Binge drinking is bad, it leads to a lot of bad things, 
and some epi stats, which means we're important.

Binge drinking interventions often try to target...something.

But those somethings are often poorly operationalized or not at all, either through...idk...probably
self-report. Studies on self-report show
(what, i imagine they're pretty bad, but maybe they're actually surprisingly good). 
Wearable transdermal alcohol sensors have emerged as an alternative to self-report,
offering a passive continuous measure of alcohol consumption
\parencite{fairbairn2025,gunn2023}. It surpasses previous measures of intoxication (eg
breathalyzers) by reducing the burden traditionally involved in obtaining a measurement to (blakn)
 what were the findings of those feasability studies?  
But although we have all this data, we dont know what to do with it.
Raw transdermal alcohol concentration (TAC) signal
can be difficult to interpret and translate into meaningful drinking parameters
\parencite{gunn2023,kianersi2023}.

To bridge the gap between theory and method, we propose a formal model for binge-drinking, which operationalizes
binged drinking by...---some number--- parameters characterizing...---the things we
characterize---. 

Drawing from previous work on binge-drinking we conceptualize
binge-drinking as consisting of discrete drinking episodes of varying magnitude
separated by a period of non-drinking \parencite{greenfield2000}.
For example, TAC-derived features such as rise rate, peak concentration, and rise duration
have been shown to predict alcohol-induced blackouts in college drinkers
\parencite{richards2024}, and TAC signal processing models have been developed and validated
against self-reported and breath alcohol concentration data in similar populations
\parencite{kianersi2023}.
Our model uses ideas from...---insert literature, ie michaelis menten---
...to characterize---the big picture formulation ie between drinking time, within episode
duration, within episode drinking patterns---as a dynamical system. 
Compared to previous models, ours provide interpretable parameters with mechanistic meanings.
Furthermore, we are able to distinguish between the end of a drinking episode and 
a return to baseline TAC. We show that the parameters can be recovered with typical
TAC sample sizes.

simulation demonstrate the accuracy of our method/software.

We illustrate how these parameters could be used to characterize the effects of risk factors 
on binge-drinking in college students.

the remainder of this paper is organized as follows: ---big long run-on sentence---

\subsection{the criteria for a good model}
We want the theory to achieve the following things:
1. It should operationalized by mechanistic parameters that are directly actionable 
by interventions.

2. It should be able to recreate things like disproportionate next day drinking 
consequences between males and females given similar rates. (ie AUC should be higher
for people with lower metabolisms). If there have been disproportionate predictive validity
in things like rise rate/fall rate/peak for predicting next day consequences based on sex,
those same discontinuities should be accounted for by the metabolism difference.

3. It should be able to recreate the TAC "tail" behavior (is it linear or exponential?).




 

\subsection{the theoretical model}

While previous work has focused on mathematical modeling of between-person population
dynamics \parencite{gruenewald1994,duffy1992}, we focus on creating a within-person
model of binge-drinking.

We want the model to capture the rate of entering an episode when in a 
non-drinking episode, the rate of exiting an episode when in a drinking episode, 
the within-episode dynamics, and the metabolism of alcohol.

Let $S(t) \in [0, 1]$ be the state of the individual at time $t$, where $S(t) = 0$ 
indicates a non-drinking episode and $S(t) = 1$ indicates a drinking episode. Our theory
deconstructs binge-drinking into four components: the rate of entering a drinking
episode, the rate of leaving a drinking episode, the within-episode ingestion rate,
and the ethanol metabolism rate.

This could be formalized by,
 \begin{equation}
 TAC(t) dt = S(t) I(t) - M(t).
 \end{equation}
Where $I(t)$ represents the alcohol ingestion rate at time $t$ and $M(t)$ represents the 
metabolic rate at time $t$. S(t) acts as a switching function, turning the intake rate 
function on or off depending on whether the individual is in a drinking episode or not.
This is similar to previous work by \parencite{duffy1992} in that we model the time between
drinking episodes as exponential and the number of drinks within each episode as Poisson
distributed. We extend this work by explicitly relating these distributions to TAC at time $t$
through a two-state continuous Markov process to characterize the between-episode dynamics 
and use of shot-noise process to characterize the within-episode dynamics.
The switching function $S(t)$ is a two-state continuous Markov process with the 
transition rates $q_{01}$ and $q_{10}$, representing the rates of entering and 
exiting a drinking episode, respectively. 

There are a few options for modeling $I(t) + M(t)$.
First is the model proposed by \parencite{giraldo2017a}, which used a dampened spring 
oscillator to model the return to equilibrium
\begin{equation}
    TAC(t) d^2t = -k^2TAC(t) - 2 \zeta k \frac{dTAC(t)}{dt} + k^2 \eta I(t),
\end{equation}
where $k > 0$ is the natural frequency of the oscillator, $\zeta > 1$ is the damping ratio, 
and $\eta > 0$ is a scaling factor for the input function $I(t)$. One straightforward
extension would be to model add a regime-switching component, so it handles the episodic
nature of drinking.

It's the Kalman prediction-error decomposition, summed over the 5 ids:

$$-2\log L=\sum_{i=1}^{n}\sum_{t=0}^{T-1}\left[\log 2\pi+\log F_{it}+\frac{v_{it}^2}{F_{it}}\right]$$

with $v_{it}=y_{it}-Hx_{it|t-1}$ and $F_{it}=HP_{it|t-1}H'+\texttt{mnoise}$. Discretisation is exact via Van Loan: $A_d=e^{A\Delta t}$, and $Q_d$ from the $\exp$ of $\begin{bmatrix}-A&Q_c\\0&A'\end{bmatrix}\Delta t$. Order per step is update-then-predict, with the $t=0$ observation included.
